\scanpage{13}
\mspage{\textcircled{1}}
\margin{1982}
\diagram{\begin{tikzcd}[column sep=large]
\mathcal X_{g,\nu\mathbb Q}\simeq M_{g,\nu+1,\nu\mathbb Q}\arrow[r,no head]\arrow[d,no head]&\mathcal X_{g,\nu\mathbb C}\simeq M_{g,\nu+1,\nu\mathbb C}\arrow[r,no head]\arrow[d,no head]&\mathcal X_{g,\nu}^{an}\simeq M_{g,\nu+1,\nu}^{an}\arrow[d,no head]\\
M_{g,\nu\mathbb Q}\arrow[r,no head]\arrow[d,no head]&M_{g,\nu\mathbb C}\arrow[r,no head]\arrow[d,no head]&M_{g,\nu}^{an}\\
\mathbb Q\arrow[r,no head]&\mathbb C&
\end{tikzcd}}

\[
\left\{\begin{array}{l}
X_{g,\nu}\ \text{courbe alg/}\mathbb C\quad\text{type }g,\nu,\ \text{corr. à}\quad
\xi_0\in\operatorname{Ob}\uline{\mathrm{Pt}}(M_{g,\nu\mathbb Q};\mathbb C)\simeq\operatorname{Ob}\uline{\mathrm{Pt}}(M_{g,\nu\mathbb C}/\mathbb C)\\
\phantom{X_{g,\nu}}\simeq(\mathcal X_{g,\nu})_{\xi_0}\hfill\simeq\operatorname{Ob}\uline{\mathrm{Pt}}(M_{g,\nu}^{an})
\end{array}\right.
\]
\[
x_0\in X_{g,\nu}(\mathbb C)\simeq X_{g,\nu}^{an},\quad (X_0,x_0)\ \text{corr. :}\quad
\xi_0\in\operatorname{Ob}\uline{\mathrm{Pt}}(\mathcal X_{g,\nu\mathbb Q};\mathbb C)\simeq\operatorname{Ob}\uline{\mathrm{Pt}}(\mathcal X_{g,\nu\mathbb C}/\mathbb C)
\deleted{\unsure{\simeq}}\simeq\operatorname{Ob}\uline{\mathrm{Pt}}(\mathcal X_{g,\nu}^{an})
\]
au dessus de $\xi_0$

\uline{Déf}
\[
(1)\quad\left\{\begin{array}{l}
\Gamma_{g,\nu}^{+}=\pi_1(M_{g,\nu}^{an},\xi_0)\\
\mathcal G_{g,\nu}^{+}=\pi_1(\mathcal X_{g,\nu}^{an},x_0)\simeq\Gamma_{g,\nu+1,\nu}\subset\Gamma_{g,\nu+1}\\
\Pi_{g,\nu}=\pi_1(X_{g,\nu},x_0)\\[4pt]
\mathcal N_{g,\nu}=\pi_1(M_{g,\nu\mathbb Q},\xi_0)\\
\mathcal M_{g,\nu}=\pi_1(\mathcal X_{g,\nu\mathbb Q},x_0)
\end{array}\right.
\]
\deleted{\illegible{1}}

d'où
\diagram{\begin{tikzcd}[column sep=small]
(2)\quad1\arrow[r]&\Pi_{g,\nu}\arrow[r]\arrow[d,equal]&\mathcal G^+_{g,\nu}\arrow[r]\arrow[d,hook]&\Gamma^+_{g,\nu}\arrow[r]\arrow[d,hook]&1\\
1\arrow[r]&\Pi_{g,\nu}\arrow[r]&\operatorname{Aut}_{loc}(\Pi_{g,\nu})\arrow[r]&\operatorname{Autext}_{loc}(\Pi_{g,\nu})&
\end{tikzcd}}

\deleted{et} et, en passant aux complétés \deleted{et}
\diagram{\begin{tikzcd}[column sep=small]
(3)\quad1\arrow[r]&\widehat\Pi_{g,\nu}\arrow[r]\arrow[d,equal]&\widehat{\mathcal G}^+_{g,\nu}\arrow[r]\arrow[d,hook,"?"']&\widehat\Gamma^+_{g,\nu}\arrow[r]\arrow[d,hook,"?"]&1\\
1\arrow[r]&\widehat\Pi_{g,\nu}\arrow[r]&\operatorname{Aut}_{loc}(\widehat\Pi_{g,\nu})\arrow[r]&\operatorname{Autext}_{loc}\unsure{(}&
\end{tikzcd}}

\deleted{$\widehat\Pi_{g,\nu}\overset{can}{\simeq}\pi_1(X_{g,\nu\mathbb C},x_0)$}
\[
(4)\quad\left\{\begin{array}{l}
\widehat\Gamma^+_{g,\nu}\overset{can}{\simeq}\pi_1(M_{g,\nu\mathbb C},\xi_0)\overset{can}{\simeq}\pi_1(M_{g,\nu\overline{\mathbb Q}},\xi_0)\\
\widehat{\mathcal G}^+_{g,\nu}\overset{can}{\simeq}\pi_1(\deleted{M}\mathcal X_{g,\nu\mathbb C},x_0)\overset{can}{\simeq}\pi_1(\mathcal X_{g,\nu\overline{\mathbb Q}},x_0)\simeq\widehat\Gamma_{g,\nu+1,\nu}\subset\widehat\Gamma_{g,\nu+1}\\
\widehat\Pi_{g,\nu}\overset{can}{\simeq}\pi_1(X_{g,\nu\mathbb C},x_0)\simeq\pi_1(X_{g,\nu\overline{\mathbb Q}},x_0)
\end{array}\right.
\]
\deleted{$\unsure{\widehat\Pi_{g,\nu}}$}

\diagram{\begin{tikzcd}[column sep=large]
&1&1&&\\
1\arrow[r]&\widehat\Gamma^+_{g,\nu}\arrow[r]\arrow[u]&\mathcal N_{g,\nu}\arrow[r,no head]\arrow[u]&\Gamma_{\overline{\mathbb Q}/\mathbb Q}\arrow[r,no head]&1\\
(5)\quad1\arrow[r,no head]&\widehat{\mathcal G}^+_{g,\nu}\arrow[r]\arrow[u]&\mathcal M_{g,\nu}\arrow[r,no head]\arrow[u]&\Gamma_{\overline{\mathbb Q}/\mathbb Q}\arrow[r,no head]\arrow[u,equal]&1\\
&\widehat\Pi_{g,\nu}\arrow[r,Rightarrow]\arrow[u]&\widehat\Pi_{g,\nu}\arrow[u]&&\\
&1\arrow[u]&1\arrow[u]&&
\end{tikzcd}}
\[\underbrace{\hspace{5em}}_{(3)}\]

\margin{\deleted{$\unsure{(6)\ 1\to\widehat\Pi_{g,\nu}\to\mathcal M_{g,\nu}\to\mathcal N_{g,\nu}\to1}$}}
\margin{\uline{NB} $\widehat\Gamma_{g,\nu},\widehat{\mathcal G}_{g,\nu}$ s'obtiennent comme images inverses de $\{1,\tau\}\subset\Gamma_{\overline{\mathbb Q}/\mathbb Q}$ ($\tau\in\Gamma_{\overline{\mathbb Q}/\mathbb Q}$, conj. complexe) dans $\mathcal N_{g,\nu},\mathcal M_{g,\nu}$}

\uline{NB} On a (via $\mathcal N_{g,\nu}\to\operatorname{Autext}_{loc}\widehat\Pi_{g,\nu}$ p.ex.)
\diagram{\begin{tikzcd}
\Gamma^+_{g,\nu}\arrow[r]&\widehat\Gamma^+_{g,\nu}\arrow[r]&\mathcal N_{g,\nu}\arrow[r,Rightarrow]&\boxed{\mathfrak S_\nu}\\
(\unsure{7})\quad\mathcal G_{g,\nu}\arrow[r,no head]\arrow[u]&\widehat{\mathcal G}_{g,\nu}\arrow[r,no head]\arrow[u]&\mathcal M_{g,\nu}\arrow[u]&
\end{tikzcd}}

d'où, en prenant les noyaux, des groupes
\[
(8)\quad\left\{\begin{array}{l}
\Gamma^{+!}_{g,\nu},\ \widehat\Gamma^{+!}_{g,\nu},\ \mathcal N^!_{g,\nu}\\
\mathcal G^{+!}_{g,\nu},\ \widehat{\mathcal G}^{+!}_{g,\nu},\ \mathcal M^!_{g,\nu}
\end{array}\right.
\]
s'interprétant comme des $\pi_1$ également ...

\mspage{12}
Il y a des des suites ``en $!$''
\deleted{\unsure{TSVP}} $(2')\ (3')\ (5')$

\figure{\begin{tikzpicture}\draw (0,0) ellipse (.65 and .25);\node at (0,0) {TSVP!};\end{tikzpicture}}
